\documentclass[10pt,a4paper]{amsart}
\usepackage[margin=19mm]{geometry}
\usepackage{amsmath,amssymb,mathtools}
\usepackage[scr=rsfs,cal=euler]{mathalfa}
\usepackage{xfrac}
\usepackage[unicode,hidelinks,bookmarks=false]{hyperref}
\newcommand{\ie}{i.e.\ }
\newcommand{\cf}{cf.\ }
\newcommand{\resp}{respectively\ }
\newcommand{\Subsection}{\subsection}
\providecommand{\nobreakdash}{\nobreak\hbox{-}}
\setlength{\emergencystretch}{2em}
\multlinegap0pt
\usepackage{bm}
\usepackage{bbm}
\usepackage{dsfont}
\usepackage{xspace}
\usepackage{cancel}
\usepackage{mathtools}
\usepackage{amsthm}
\makeatletter
\@ifundefined{theorem}{\newtheorem{theorem}{Theorem}}{}
\makeatother
\DeclareMathOperator{\Z}{\mathbb{Z}}
\title[Unique continuation on planar graphs]{Unique continuation on planar graphs}
\author{Ahmed Bou-Rabee, William Cooperman, Shirshendu Ganguly}
\date{Source: arXiv:2309.13728v2 (CC BY 4.0)}
\begin{document}
\maketitle

\noindent\emph{Source and licence.} Unique continuation on planar graphs. Version arXiv:2309.13728v2 (CC BY 4.0), distributed under the Creative Commons Attribution 4.0 International licence, as recorded in the arXiv licence metadata for this exact version and shown on the abstract page https://arxiv.org/abs/2309.13728v2. Full rights notice: licenses/arxiv-2309.13728-CC-BY-4.0.txt.\par\smallskip

\noindent\textit{Editorial scope.} Counted unit: the Theorem whose source label is \texttt{theorem:counterexample} in the article (source0.tex lines 381--386, source label \texttt{theorem:counterexample}), delivered together with the article's own complete original proof of it (source0.tex lines 388--413, one \texttt{proof} environment of 26 lines). No mathematics is written, repaired or re-derived: statement and proof are the article's own text line for line. Required context retained uncounted: none. Every definition, display and label of those lines is retained unchanged. Omitted from this package and not redistributed: 1--380, 387--387, 414--560. Nothing inside a retained excerpt is omitted or altered: every retained body is delivered exactly as the article's lines are written, so no omission receipt of any kind accompanies this package. Citations: the retained excerpts carry no citation command at all, so no bibliography entry of the article is transcribed and no reference is redistributed. Reference adaptation: none was needed and none was made; every reference inside the delivered unit resolves inside it, so no unresolved reference or citation is produced. Printed slips kept literal: no printed word, symbol, hypothesis or number of the counted statement or of its proof was altered. Layout and class: the article's own class is \texttt{daj} with packages including hyperref, graphicx, amsmath, amsfonts, amssymb, amsthm, mathrsfs; the wrapper is the lane's pinned \texttt{amsart} editorial wrapper at 10pt on A4, which restates the theorem-like environments the retained text needs and adds the article's own macro definitions that the retained text needs (\texttt{\textbackslash Z}), with extra wrapper packages (bm, bbm, dsfont, xspace, cancel, mathtools). The delivered numbering is the unit's own. Assets: no retained span contains a figure environment, a graphics-inclusion command, a PDF-page inclusion, a table or a TikZ picture; no image file is copied into this package, the package declares no asset dependency and no image file is redistributed with it. Licence: version arXiv:2309.13728v2 of this article is distributed under the Creative Commons Attribution 4.0 International licence, as recorded in the arXiv licence metadata for this exact version and shown on the abstract page https://arxiv.org/abs/2309.13728v2. A full rights notice is delivered at licenses/arxiv-2309.13728-CC-BY-4.0.txt. Characters: every retained excerpt is pure ASCII, so the delivered engine, XeLaTeX with its default Latin Modern text font, reports no missing character.
\begin{theorem}\label{theorem:counterexample}
	For each choice of positive $A_1 \neq A_2$
	and $A_3 > \frac{2 A_1^2 A_2^2}{{(A_1-A_2)}^2 (A_1+A_2)}$
	there is a choice of $A_4 > 0$ for which there is a harmonic
	function on $(\Z^2, \mathcal{E}(A_1,A_2,A_3,A_4))$ supported on the diagonal line $\{x_1 = x_2\}$.
\end{theorem}

\begin{proof}
	For the sequence $\{z_i\}_{i \in \Z}$ determined by the equalities $z_0 = 1$ and $A_1 z_{i}  + A_2 z_{i-1} = 0$,
	a short computation shows that
	the function $h(x_1, x_2) = 1_{\{ x_1 = x_2\}} z_{x_1}$
	is harmonic on $(\Z^2, \mathcal{E}(A_1,A_2,A_3,A_4))$ for some $A_4 > 0$.
	% By evaluating the equation at the off-diagonal, we have
	% \[
	%   A z_{i}  + A_2 z_{i-1} = 0 \,.
	% \]
	% Thus, after fixing $z_0 = 1$, the above equation determines $\{z_i\}_{i \in \Z}$. Next, by evaluating the equation on the diagonal $x_1 = x_2$,  we have
	% \[
	%   2 (A + A_2 + A_3 + A_4) z_{x_1} =  A_4 (z_{x_1+1} + z_{x_1-1}) + A_3 ( z_{x_1-2} + z_{x_1+2}) \, .
	% \]
	% Hence, by plugging in the values of $z_x$, we have
	% \[
	%   2(A + A_2 + A_3 + A_4) = - A_4 \left( \frac{A}{A_2}  + \frac{A_2}{A} \right)  + A_3\left( \frac{A_2^2}{A^2} + \frac{A^2}{A_2^2} \right) \, .
	% \]
	% From this equation, we see that for every
	% choice of $A_3$ satisfying the inequality in the statement,
	% $A_4$ is positive.
	% Indeed, we see that $A_4 = (A^3 A_3 - A^2 A_2 ( 2A_2 + A_3) - A A_2^2 A_3 + A_2^3 A_3) (A A_2( A + A_2))^{-1}$.
	% Thus, for every
	% choice of $A_3$ satisfying the inequality in the statement,
	% $A_4$ is positive.
	% Since, at all other vertices, harmonicity gives no constraints, \ie, we have the equality $0 = 0$, this completes the proof.
\end{proof}

\end{document}
